\section{設定と主定理}\label{sec:main}

全ての多様体、形式、写像は、特に断らない限り滑らかである。
Hamiltonian vector fieldは$\iota_{X_H}\omega=-\mathrm dH$で定義し、normal circleの周期は$2\pi$とする。
従って標準複素直線上の正のmomentは$|z|^2/2$である。
$\dd^c$は$\dd^c\log|z|^2=\dd\theta$となる規約を用いる。
$H$はhyperplane classで、$\int_{\PP^1}H=1$と正規化する。

\begin{definition}\label{def:fibration}
本稿の\emph{分岐底空間上のLagrangian torus fibration}とは、コンパクトなsymplectic 4多様体$(X,\omega)$から有限2複体$B$への、連結なfibersを持つproper continuous surjection $\pi:X\to B$で、$B$の稠密な開2次元stratum上では滑らかなsubmersionとなり、そのfibersがLagrangian $T^2$となるものをいう。
特異fiberには、以下に指定する局所的な層構造を与え、その滑らかな各stratumはisotropicとする。
ここでstratificationは\cite[Definition 2.1]{EM}の弱い意味、すなわち閉部分集合のfiltrationの各差が滑らかな多様体となるという意味で用い、Whitney条件は含めない。
\emph{自由な辺がない}とは、2次元stratumの閉包に接する各開1次元stratumに、ちょうど一枚だけのsheetが接することがない、という意味である。
\end{definition}

この定義は、基底が通常の2次元多様体であるという条件を含まない。
本稿では分岐辺の横断的な断面は三叉グラフ$Y$である。
正則な底空間にはaction-angle theoremによるintegral affine structureがあるが、分岐を横断する単一のaffine chartは定義からは得られない。

各$i\ne j$に対し、$d-2$個の葉が$\{1,\ldots,d\}\setminus\{i,j\}$でラベル付けされた三価木$T_{ij}=T_{ji}$を選ぶ。
これは実際の多穴球面$C_{ij}^{\circ}$上のmarked pants分解で実現し、単なる抽象木の選択とは区別する。
$B_d(T)$は、各成分のcoreを一つの頂点に、各$T_{ij}$の内部頂点を$i,j$間の分岐辺に、内部辺を二角形に対応させ、各三重点を三角形に対応させた2複体である。
厳密なcellular定義と接着写像は底空間の節で与える。

\begin{theorem}[境界のない分岐底空間への構成]\label{main:global}
$X_d\subset\PP^3$を滑らかな次数$d\ge5$の曲面、$\omega$を$[\omega]=aH$, $a>0$を満たすKähler formとする。
任意に選んだ上記のmarked pants treesを用いて
\[
\pi_d:(X_d,\omega)\longrightarrow B_d(T)
\]
という定義\ref{def:fibration}のfibrationを構成できる。そのfiberは次の表の通りであり、他の型の特異fiberを生じない。
\begin{center}
\begin{tabularx}{\linewidth}{@{}lXr@{}}
\toprule
底空間のstratum & fiber & 個数または成分数\\
\midrule
正則な開2次元部分 & Lagrangian $T^2$ & ---\\
開分岐辺 & $S^1\times\figeight$ & $\binom d2(d-4)$\\
面内のnodal values & 一度pinchされた$T^2$ & $d\binom d2$\\
core頂点 & $\PP^2$の$d-1$直線補集合のlogarithmic Weinstein skeletonのsymplectic image & $d$\\
\bottomrule
\end{tabularx}
\end{center}
各分岐辺にはちょうど三枚のsheetが接する。
基底は$\binom d3$枚の三角形、$\binom d2(d-5)$枚の二角形を持ち、
\[
B_d(T)\simeq\Sk_2\Delta^{d-1}\simeq\bigvee^{\binom{d-1}{3}}S^2.
\]
$d=5$では$B_5\cong\Sk_2\Delta^4$であり、$d\ge6$では$B_d(T)$は$\Sk_2\Delta^{d-1}$と同相ではない。
\end{theorem}

coreの同定は、適切に選んだlogarithmic Weinstein構造のskeletonとの同定である。
特定のembedded skeletonを最初から指定して保存する、という主張ではない。
Morse型coreの等方的なstratificationには\cite[\S2.3]{EM}を用いる。
core fiberのEuler数は$\binom{d-3}{2}$であり、$d$とともに変化する。
従ってこの定理は、全ての次数に共通の一種類のpants fiberだけで構成するという結論も含まない。

\begin{theorem}[vertical Kähler slice]\label{main:cone}
後述の順序付き小解消$f:\widetilde{\cX}_d\to\Delta$の中心成分を$D_1,\ldots,D_d$とする。
$a>0$を固定し、$\sum_i[D_i]=0$による共通定数の不定性を除いて
\[
\xi=aH+\sum_{i=1}^d\beta_i[D_i]
\]
を考える。
底を十分縮小した全空間上のKähler classとして$\xi$が実現できるための必要十分条件は
\[
\beta_1<\beta_2<\cdots<\beta_d,
\qquad a+\sum_i\beta_i-d\beta_d>0
\]
である。$\delta_k=\beta_{k+1}-\beta_k$を用いると、この領域は
\[
\delta_k>0\quad(1\le k\le d-1),\qquad
\sum_{k=1}^{d-1}k\delta_k<a
\]
という開単体になる。
\end{theorem}

これは全相対Kähler coneの記述ではなく、指定したvertical affine sliceとの交わりの完全な記述である。
大域fibrationの証明には、この単体の中の十分小さいrayを用いる。
単体の全ての点で同じ局所surgeryが実行できると主張する必要はない。

\paragraph{証明の構造.}
まず明示的な平面退化と偏極を構成する。
次に、直線配置補集合の有限体積compactificationを、境界上のmarked Morse葉層に相対的に作る。
小さいdivisor-compatible blow-upによってfocus--focus fibersを挿入し、実際の中心成分へfibrationを移す。
その後、divisorから正の距離を置いた飽和領域を、同じ全空間のsymplectic connectionで平滑化へ運ぶ。
二重曲線ではmomentに依存するMorse葉層を補間し、三重点ではtoric moment cutsを使う。
この順序により、nodeのコンパクトなfiber全体を保ち、gluingの際にfiberを途中で切ることを避ける。

\paragraph{一般化と局所的な障害.}
上の平面模型に必要な相対 compactification は，正の重み付き Kähler divisor 表示を持つ
任意の genus の SNC divisor に対して構成する（定理\ref{snc:main}）．
これにより，$d=e_1+\cdots+e_k$，$k\geq3$ の次数分割に対し，
$k$ 個の core と $d\sum_{i<j}e_i e_j$ 個の nodes を持つ構成を得る
（定理\ref{part:global}）．その底空間は
\[
 B_{\boldsymbol e}\simeq
 \bigvee^{\,p_g(X_d)-\sum_i p_g(S_i)}S^2
\]
であり，整数鎖の強変形収縮を明示する（定理\ref{part:homotopy}）．
ここで $S_i\subset\PP^3$ は次数 $e_i$ の smooth surface である．

この構成の core を保つ場合，三価 Bergman link の分類から $d\geq7$ の全ての次数分割に
局所位相障害が生じる（定理~\ref{trop:cubic-classification}，命題~\ref{trop:partition-obstruction}）．
六次では全平面分割だけが残り，六つの同時 Petersen systems がある．
五次では全平面と $(1,1,1,2)$ だけが全 core の matroidal local topology を持ち得る．
後者の標識付き抽象 graph systems は四個である（定理~\ref{q5:classification}）．

\paragraph{作用適合 atlas の存在と分類.}
作用適合 matroidal atlas の transitions は，ambient integer affine maps の制限とする．
正則 faces の余接線格子は actual fibration の作用格子に等しい．
focus points を通常の matroidal charts から除き，その punctured neighborhoods に
実際の focus affine singularities を残す（定義~\ref{atl:definition}）．
branch denominator $q$ は，primitive branch tangent に対する正 normal action の
微分の絶対値であり，以下の構成は全 branch で $q=1$ を満たす．

\begin{theorem}[五次・六次の大域作用適合 atlas]\label{main:action-atlas}
$d=5,6$，$X_d\subset\PP^3$ smooth とし，$\omega$ を
$[\omega]=aH$，$a>0$ を満たす Kähler form とする．
定理~\ref{main:global} の全平面 fibration を，focus affine singularities を許す
大域作用適合 matroidal atlas を持つように選べる．
core charts は五次では標準 $K_4$，六次では標準 Petersen fans である．
五次の十 triangles には各五つ，六次の二十 triangles と十五 bigons にはそれぞれ三つと二つの
nodal values が入る．全 branch の正規化された affine length は $a/(\pi(d-1))$ である．
\end{theorem}

\begin{remark}[作用 atlas と KE 計量実現の選択]\label{main:atlas-ke-scope}
定理\ref{main:action-atlas}は，指定された Kähler form ごとに fibration を選ぶ存在定理である．
各小さい $t\ne0$ の actual 負 KE form もこの対象となる．
五次ではさらに，系\ref{ke:quintic-simultaneous-main}と
定理\ref{ke:quintic-simultaneous-realization}により，作用 atlas を持つ同じ射影を
全曲面 ambient 距離の GH 実現にも使える．等重みの対称四角形系を選べば，
各 triangle に五つの nodal values が入る．
その証明では actual primitive の両端切片と整数 framing を有限式で合わせ，
初期 source 葉の全 tail を小面積に選び，source wall を含む全 moment の面積輸送を行う．
同じ canonical moment，保持 corner の全葉，中央 nodal neighborhoods と
全 overlap の whole-fiber partition を保つため，
命題\ref{ke:moment-relative-replacement}から同じ射影の距離比較を得る．
任意に先指定した source fibers・大きい保持 annuli・critical values への適合は要求しない．
六次でも，定理\ref{ke:sextic-simultaneous-realization}が標準等重み Petersen 系について
同じ射影による atlas・GH 実現・確率体積極限を与える．
その追加内容は，内部二 source block の累積位置を含む二端 cut，
各側三 sources の separator による全 moment 実現，中央 source 全近傍を保つ
parameter の固定，及び平行移動まで一致する三面の整数アフィン接続である．
\end{remark}

定理\ref{main:action-atlas}の証明は，局所作用座標，整数 framing，critical-area を指定した actual neck 接着を
まとめる定理~\ref{ga8:global} で与える．
個別の core charts を並べるだけでは，作用の加法定数と atlas transitions は決まらない．
本稿は固定 nonzero moments 上の exact symplectic transport と相対 Moser により
この接着条件を満たす．

混合分割 $(1,1,1,2)$ では，三つの $K_4$ と一つの $K_{3,3}$ の各 core に個別の
作用適合 chart がある（系~\ref{mixper:individual}）．
しかし，四 puncture neck の両端の annulus を比較すると，正の分母 $q_A,q_B$ は
$q_Aq_B\mid2$ を満たし，三価 branch の奇数性により両方とも一となる
（命題~\ref{q9:rigidity}）．この剛性と法円作用の分岐係数の総和 $-6$ により，
混合五次の大域接続は任意の分母を許しても不可能である
（定理~\ref{hd9:mixed-no-go}）．

\begin{theorem}[現行 SNC 構成の matroidal atlas の分類]\label{main:atlas-classification}
$d\ge5$，$d=\sum_{i=1}^ke_i$，$k\ge3$ とする．本稿の比例曲率 SNC cores，
既存の branch graph，標準 coarse matroidal cones を保ち，追加の branch singularities を
入れない範囲で，大域作用適合 matroidal atlas を持つ fibration を構成できる次数分割は，
全平面五次と全平面六次に限る．branch denominators に一という仮定を課さない．
\end{theorem}

証明は系~\ref{hd9:classification}で与える．全平面六次の atlas では，
三十本全ての branch denominators が自動的に一となり，各 bigon は二 nodes を持つ
（系~\ref{q9:sextic-rigidity}，注意~\ref{q9:node-budget}）．
一方，四直線の一つの core では任意の奇数 $q\ge3$ に対し分母 $(q,1,1,1)$ の chart が
存在する（\cite{QuinticArithmetic}）．従って分母剛性は大域接着を使う条件である．
分類の範囲は指定した SNC 構成であり，core surgery，fan の変更，追加の branch singularities
を許す全ての fibrations を分類する主張ではない．

\paragraph{五次の整数アフィン接着との関係．}
全平面五次で任意の奇数分母を許す chart の同時接着については，局所の整数行列，
共有 puncture の整数 clutching，node 配分，および対向する分岐係数の相殺を同時に調べる必要がある．
これらの算術条件，高分母支持 graph に対する障害，および局所整数族の分類は，
companion 論文\cite{QuinticArithmetic}にまとめる．
定理~\ref{main:action-atlas}で構成する分母一の大域 atlas の存在証明は，本稿内で完結する．
